\documentclass[../../../include/open-logic-section]{subfiles}

\begin{document}

\olfileid{sth}{spine}{Valphabasic}
\olsection{તબક્કાઓના મૂળભૂત ગુણધર્મો}

ગણસિદ્ધાંતના પાયા માટે $V_\alpha$ની વ્યાખ્યાનું મહત્ત્વ
સ્પષ્ટ કરવા તેના કેટલાક મૂળભૂત પરિણામો આપીશું. શરૂઆત એક
વ્યાખ્યાથી કરીએ:\footnote{આ ગુણધર્મ માટે કોઈ પ્રમાણભૂત નામ નથી;
અહીંનું નામ \citet{ButtonLT1}ના નામનો અર્થ સ્પષ્ટ કરે છે.}
\begin{defn}
	ગણ $A$ \emph{ઉપગણ-સંવૃત} છે જો અને માત્ર જો
	$\forall x((\exists y \in A)x \subseteq y \lif x \in A)$.
\end{defn}
\begin{lem}\ollabel{Valphabasicprops}
દરેક ક્રમવાચક સંખ્યા $\alpha$ માટે:
\begin{enumerate}
	\item\ollabel{Valphatrans} દરેક $V_\alpha$ પરંપરિત છે.
	\item\ollabel{Valphapotent} દરેક $V_\alpha$ ઉપગણ-સંવૃત છે.
	\item\ollabel{Valphacum} જો $\gamma \in \alpha$, તો $V_\gamma
	\in V_\alpha$ (અને તેથી \olref{Valphatrans} મુજબ
	$V_\gamma \subseteq V_\alpha$ પણ).
\end{enumerate}
\end{lem}

\begin{proof}
આ પરિણામો (એકસાથે) અનંતાતીત અનુમાનથી સાબિત કરીએ. અનુમાન માટે
ધારો કે દરેક ક્રમવાચક સંખ્યા $\beta < \alpha$ માટે
\olref{Valphatrans}--\olref{Valphacum} સાચાં છે.

$\alpha = \emptyset$નો કિસ્સો સીધો છે.

ધારો કે $\alpha = \ordsucc{\beta}$. \olref{Valphacum} માટે,
જો $\gamma \in \alpha$ હોય, તો અનુમાનની ધારણા મુજબ
$V_\gamma \subseteq V_\beta$, તેથી
$V_\gamma \in \Pow{V_\beta} = V_\alpha$. \olref{Valphapotent}
માટે ધારો કે $A \subseteq B \in V_\alpha$, એટલે કે $A
\subseteq B \subseteq V_\beta$; ત્યારે $A \subseteq V_\beta$,
તેથી $A \in
V_\alpha$. \olref{Valphatrans} માટે ધ્યાન આપો કે $x \in A \in
V_\alpha$ હોય તો $A \subseteq V_\beta$, એટલે $x \in V_\beta$.
અનુમાનની ધારણા મુજબ $V_\beta$ પરંપરિત હોવાથી $x
\subseteq V_\beta$, અને તેથી $x
\in V_\alpha$.

ધારો કે $\alpha$ સીમા ક્રમવાચક સંખ્યા છે. \olref{Valphacum}
માટે જો $\gamma \in \alpha$ હોય, તો
$\gamma \in \ordsucc{\gamma} \in \alpha$; તેથી અનુમાનની ધારણા
મુજબ $V_\gamma \in V_{\ordsucc{\gamma}}$, અને તેથી
$V_\gamma \in \bigcup_{\beta \in \alpha} V_\beta = V_\alpha$.
\olref{Valphatrans} અને \olref{Valphapotent} માટે માત્ર
એટલું ધ્યાનમાં લો કે પરંપરિત ગણોનો યોગ પરંપરિત છે અને
ઉપગણ-સંવૃત ગણોનો યોગ ઉપગણ-સંવૃત છે.
\end{proof}

\begin{lem}\ollabel{Valphanotref}
દરેક ક્રમવાચક સંખ્યા $\alpha$ માટે $V_\alpha \notin V_\alpha$.
\end{lem}

\begin{proof}
અનંતાતીત અનુમાનથી સાબિત કરીએ. સ્પષ્ટ છે કે
$V_\emptyset \notin V_\emptyset$.

જો $V_{\ordsucc{\alpha}} \in V_{\ordsucc{\alpha}} = \Pow{V_\alpha}$,
તો $V_{\ordsucc{\alpha}} \subseteq V_\alpha$; અને
\olref{Valphabasicprops} મુજબ $V_\alpha
\in V_{\ordsucc{\alpha}}$ હોવાથી
$V_\alpha \in V_\alpha$. આમ તેનું પ્રતિવિધાન મળે છે: જો
$V_\alpha \notin V_\alpha$ હોય, તો
$V_{\ordsucc{\alpha}} \notin V_{\ordsucc{\alpha}}$.

જો $\alpha$ સીમા ક્રમવાચક સંખ્યા હોય અને
$V_\alpha \in V_\alpha = \bigcup_{\beta \in
\alpha}V_\beta$, તો કોઈ $\beta \in
\alpha$ માટે $V_\alpha \in V_\beta$; પણ ત્યારે
$V_\beta \in V_\alpha$ પણ છે, જેથી
\olref{Valphabasicprops}નો બે વાર ઉપયોગ કરતાં $V_\beta \in
V_\beta$ મળે છે. આમ પ્રતિવિધાન પ્રમાણે, જો બધા $\beta \in \alpha$
માટે $V_\beta
\notin V_\beta$ હોય, તો $V_\alpha \notin
V_\alpha$.
\end{proof}

\begin{cor}
કોઈ પણ ક્રમવાચક સંખ્યાઓ $\alpha, \beta$ માટે:
$\alpha \in \beta$ જો અને માત્ર જો $V_\alpha \in V_\beta$.
\end{cor}

\begin{proof}
\Olref{Valphabasicprops}થી એક દિશા મળે છે. ઊલટું, ધારો કે
$V_\alpha \in V_\beta$. તો \olref{Valphanotref} મુજબ
$\alpha \neq \beta$; અને $\beta \notin \alpha$, કારણ કે
નહિતર $V_\beta \in V_\alpha$ થાય અને તેથી
\olref{Valphabasicprops}ને બે વાર વાપરતાં
$V_\beta \in V_\beta$ મળે, જે \olref{Valphanotref}નો
વિરોધ કરે છે. તેથી ત્રિવિકલ્પતા મુજબ $\alpha \in \beta$.
\end{proof}
\noindent
આ બધાથી દરેક $V_\alpha$ને પદાનુક્રમનો $\alpha$મો તબક્કો
ગણી શકીએ. કારણ આ છે.

આપણા $V_\alpha$ \emph{પુનરાવર્તી} પ્રક્રિયામાં રચાતા હોવાનું
વિચારી શકાય છે, કારણ કે ક્રમવાચક સંખ્યાઓનો આપણો ઉપયોગ
પુનરાવર્તનનો ખ્યાલ દર્શાવે છે. વધુમાં, એક તબક્કો બીજા પહેલાં
રચાતો હોય, એટલે કે $V_\beta \in V_\alpha$, એટલે કે
$\beta \in \alpha$, તો રચનાની પ્રક્રિયા \emph{સંચયી} છે,
કારણ કે $V_\beta \subseteq V_\alpha$. છેલ્લે, કોઈ પણ અગાઉના
તબક્કે ઉપલબ્ધ ગણોના \emph{બધા} સંભવિત સંગ્રહો ખરેખર રચીએ
છીએ, કારણ કે દરેક અનુગામી તબક્કો $V_{\ordsucc{\alpha}}$
એ તેના પુરોગામી $V_\alpha$નો ઘાતગણ છે.

$\ZFminus$માં ગણોના સંચયી-પુનરાવર્તી પદાનુક્રમની આપણી
કલ્પનાને સ્પષ્ટ કરવામાં આપણે \emph{લગભગ} સફળ થયા છીએ.
(જોકે પ્રતિસ્થાપનને વાજબી ઠેરવવાનું હજી બાકી છે.)

\end{document}
